\documentclass[11pt,a4paper]{article}

% Required packages
\usepackage{graphicx}
\usepackage{float}
\usepackage{caption}
\usepackage{adjustbox}

\captionsetup{
    font=small,
    labelfont=bf
}

\usepackage[margin=1.5cm]{geometry}
\usepackage{pdflscape}
\usepackage{amsmath}
\usepackage{amsfonts}

\title{Chapter 5: Architecture and Process Model\\of the RS-DRL-Based Self-Adaptive System\\(Alternative 3 - Decomposed Figures)}
\author{Thesis Documentation}
\date{\today}

\begin{document}

\maketitle

\tableofcontents
\newpage

% Set chapter counter to 5 so sections are numbered 5.1, 5.2, etc.
\setcounter{section}{0}
\renewcommand{\thesection}{5.\arabic{section}}
\renewcommand{\thesubsection}{5.\arabic{section}.\arabic{subsection}}
\renewcommand{\thesubsubsection}{5.\arabic{section}.\arabic{subsection}.\arabic{subsubsection}}

% ============================================================================
% NOTE: This version uses decomposed figures for better cognitive clarity
% Monolithic process-flow diagrams are moved to Appendix A
% Each section has focused, single-responsibility figures
% ============================================================================

% ============================================================================
% SECTION 5.1: Chapter Overview and Architectural Perspective
% ============================================================================

\section{Chapter Overview and Architectural Perspective}
\label{sec:overview}

\subsection{Chapter Scope and Objectives}

This chapter presents the architecture and process model of the RS-DRL-based self-adaptive system, focusing on system-level organization, runtime control flow, and component interaction patterns. The chapter emphasizes architectural responsibility and design rationale, maintaining clear separation between system architecture (presented here) and algorithmic mechanisms (detailed in appendices).

\subsection{Process-Oriented Organization}

The primary organizing principle is \textbf{process progression}---each section answers: ``What happens next in the system, and why?'' MAPE-K serves as an \textbf{explanatory lens}, not a competing organizational structure. This chapter uses \textbf{decomposed figures} aligned with section structure, where each figure focuses on one conceptual responsibility. Complete end-to-end workflows are provided in Appendix A for readers requiring full process context.

\subsection{Architectural Baseline}

The architectural baseline is the RS-DRL integration architecture (Figure~4 of the original article), shown in Figure~\ref{fig:rs-drl-architecture}.

\begin{figure}[H]
    \centering
    \adjustbox{width=0.95\textwidth,height=0.85\textheight,center,keepaspectratio}{\includegraphics{image/RS-DRL.png}}
    \caption{RS-DRL integration architecture showing the integration of RS-DRL with MAPE-K framework.}
    \label{fig:rs-drl-architecture}
\end{figure}

% ============================================================================
% SECTION 5.2: System-Level Self-Adaptation Loop
% ============================================================================

\section{System-Level Self-Adaptation Loop}
\label{sec:system-loop}

\subsection{Managed and Managing System Abstraction}

The RS-DRL-based self-adaptive system consists of two primary components: the \textbf{Managed System} (executes application logic) and the \textbf{Managing System} (implements adaptation logic using MAPE-K enhanced with RS-DRL).

\subsection{Continuous Adaptation Cycle (Stages 1--9)}

The system operates through a nine-stage adaptation flow:

\begin{enumerate}
    \item \textbf{Stage 1}: Runtime data collection
    \item \textbf{Stage 2}: State forwarding to RS-DRL module
    \item \textbf{Stage 3}: Asynchronous model training and storage
    \item \textbf{Stage 4}: State forwarding to Analyzer
    \item \textbf{Stage 5}: Adaptation triggering
    \item \textbf{Stage 6}: Trained model retrieval
    \item \textbf{Stage 7}: Adaptation option prediction
    \item \textbf{Stage 8}: Adaptation option retrieval
    \item \textbf{Stage 9}: Adaptation application
\end{enumerate}

All subsequent sections refine \textbf{subsets of Stages 1--9}, progressively decomposing the system-level loop into focused subprocesses.

\begin{figure}[H]
    \centering
    \adjustbox{width=0.95\textwidth,height=0.85\textheight,center,keepaspectratio}{\includegraphics{figures/rs-drl-integration-overview.pdf}}
    \caption{System-level self-adaptation process overview showing the nine-stage adaptation cycle. 
    This high-level view is decomposed into focused subprocesses in subsequent sections.}
    \label{fig:integration-overview}
\end{figure}

\subsection{Separation of Runtime Adaptation and Asynchronous Learning}

A critical architectural design decision is the temporal separation between runtime adaptation (Stages 1, 4--9) and asynchronous learning (Stages 2--3). This separation ensures runtime stability while enabling continuous policy improvement.

% ============================================================================
% SECTION 5.3: Mapping to MAPE-K Reference Model
% ============================================================================

\section{Mapping the Proposed Architecture to the MAPE-K Model}
\label{sec:mape-k-mapping}

\subsection{Conceptual Mapping to MAPE-K}

This section provides a \textbf{conceptual mapping} of the RS-DRL architecture to MAPE-K. This mapping serves as a \textbf{classification lens}, categorizing the stages according to MAPE-K functional roles.

The mapping is:
\begin{itemize}
    \item \textbf{Monitor} (Stages 1, 2, 4): Observes managed system and constructs state representations
    \item \textbf{Analyze} (Stages 4, 5): Evaluates state against goals and triggers adaptation
    \item \textbf{Plan} (Stages 5, 6, 7): Selects adaptation options using RS-DRL policies
    \item \textbf{Execute} (Stages 8, 9): Applies selected adaptations
    \item \textbf{Knowledge} (Stages 3, 6, 8): Stores models, options, and state history
\end{itemize}

\subsection{Role of RS-DRL within MAPE-K}

RS-DRL is positioned as an enhancement to the \textbf{Plan} phase, providing learned intelligence for adaptation option selection. The RS-DRL module is an \textbf{internal decision service within the Planner}, not the system architecture itself.

\subsection{Temporal Classification}

Processes are classified by temporal characteristics:

\begin{table}[H]
\centering
\begin{tabular}{ll}
\textbf{Process} & \textbf{Temporal Nature} \\
\hline
Monitoring & Continuous \\
Analysis & Event-driven \\
Planning & Event-driven \\
Execution & Event-driven \\
Learning & Asynchronous \\
\end{tabular}
\caption{Temporal classification of architectural processes.}
\label{tab:temporal-classification}
\end{table}

\begin{figure}[H]
    \centering
    \adjustbox{width=0.9\textwidth,height=0.85\textheight,center,keepaspectratio}{\includegraphics{figures/mape-k-loop.pdf}}
    \caption{MAPE-K validation view showing RS-DRL enhancement of the planning phase.}
    \label{fig:mape-k}
\end{figure}

% ============================================================================
% SECTION 5.4: Runtime Monitoring and State Construction
% ============================================================================

\section{Runtime Monitoring and State Construction}
\label{sec:monitoring}

This section describes the monitoring process using \textbf{decomposed figures} that align with subsection structure. The complete end-to-end monitoring workflow is provided in Appendix A.1.

\subsection{Runtime Data Collection}

The monitoring process begins with runtime data collection from the managed system and environment (Stage 1).

% NOTE: This would ideally use a decomposed figure showing only data acquisition
% Figure to be created: monitoring-data-acquisition.pdf (decomposed from monitoring-state-construction.pdf)
% Shows: Raw data sources → Sensor Data component → Initial collection
% Excludes: Preprocessing, feature extraction, validation (shown in subsequent figures)

\begin{figure}[H]
    \centering
    \adjustbox{width=0.75\textwidth,height=0.5\textheight,center,keepaspectratio}{\includegraphics{figures/monitoring-data-acquisition.pdf}}
    \caption{Runtime data acquisition from managed system. 
    This focused view shows data collection only (Stage 1). 
    Complete monitoring workflow in Appendix A.1.}
    \label{fig:data-acquisition}
\end{figure}

\subsection{Data Preprocessing and Feature Extraction}

Raw data is transformed through preprocessing (filtering, normalization) and feature extraction, deriving goal-relevant attributes while reducing dimensionality.

% NOTE: This would use a second decomposed figure showing only preprocessing/feature extraction
% Figure to be created: monitoring-feature-extraction.pdf
% Shows: Preprocessed Data → Feature Extractor → Feature categories
% Excludes: Raw acquisition (previous figure), validation (next figure)

\begin{figure}[H]
    \centering
    \adjustbox{width=0.75\textwidth,height=0.5\textheight,center,keepaspectratio}{\includegraphics{figures/monitoring-feature-extraction.pdf}}
    \caption{Data preprocessing and feature extraction transforming raw data into feature vectors. 
    Complete monitoring workflow in Appendix A.1.}
    \label{fig:feature-extraction}
\end{figure}

\subsection{State Vector Construction and Validation}

Extracted features are assembled into a state vector $s_t \in \mathbb{R}^d$ that serves as the interface between monitoring and decision-making. The state vector is validated to ensure consistency.

% NOTE: This would use a third decomposed figure showing only state assembly and validation
% Figure to be created: monitoring-state-validation.pdf
% Shows: Features → State Vector Assembly → Validation → Validated State Vector
% Excludes: Acquisition and preprocessing (previous figures)

\begin{figure}[H]
    \centering
    \adjustbox{width=0.75\textwidth,height=0.5\textheight,center,keepaspectratio}{\includegraphics{figures/monitoring-state-dissemination.pdf}}
    \caption{State vector construction, validation, and dissemination to downstream components. 
    Complete monitoring workflow in Appendix A.1.}
    \label{fig:state-validation}
\end{figure}

\subsection{Process Boundary}

\textbf{Process boundary}: The monitoring process ends once a validated state vector exists and has been disseminated. Everything that follows operates on this validated state representation.

For the complete end-to-end monitoring and state construction workflow, see Appendix A.1.

% ============================================================================
% SECTION 5.5: Runtime Analysis and Adaptation Triggering
% ============================================================================

\section{Runtime Analysis and Adaptation Triggering}
\label{sec:analysis}

\subsection{Goal Evaluation}

The Analyzer receives state vectors from the Monitor (Stage 4) and evaluates whether the current system state satisfies adaptation goals.

\subsection{Adaptation Triggering}

When the Analyzer detects that adaptation goals are violated or at risk, it triggers the planning phase (Stage 5).

\begin{figure}[H]
    \centering
    \adjustbox{width=0.7\textwidth,height=0.5\textheight,center,keepaspectratio}{\includegraphics{figures/mape-k-analysis-phase.pdf}}
    \caption{Analysis and adaptation triggering (Stages 4--5). 
    This focused view shows goal evaluation and triggering only.}
    \label{fig:analysis}
\end{figure}

% ============================================================================
% SECTION 5.6: Runtime Planning and Adaptation Option Prediction
% ============================================================================

\section{Runtime Planning and Adaptation Option Prediction}
\label{sec:planning}

This section describes the planning process using decomposed figures. The complete adaptation decision workflow is provided in Appendix A.2.

\subsection{Model Retrieval}

Upon receiving the adaptation trigger (Stage 5), the planning process retrieves the trained RS-DRL model from the Knowledge repository (Stage 6).

% NOTE: Decomposed figure showing only model retrieval
% Figure to be created: prediction-model-retrieval.pdf
% Shows: Adaptation trigger → Knowledge access → Model retrieval
% Excludes: Q-value evaluation, option selection (next figures)

\begin{figure}[H]
    \centering
    \adjustbox{width=0.65\textwidth,height=0.4\textheight,center,keepaspectratio}{\includegraphics{figures/prediction-model-retrieval.pdf}}
    \caption{Trained model retrieval from Knowledge repository (Stage 6). 
    Complete decision workflow in Appendix A.2.}
    \label{fig:model-retrieval}
\end{figure}

\subsection{Adaptation Option Evaluation}

The Adaptation Option Predictor evaluates Q-values for all available adaptation options using the retrieved model.

% NOTE: Decomposed figure showing only Q-value evaluation
% Figure to be created: prediction-option-evaluation.pdf
% Shows: State vector + Model → Q-value computation → Ranked options
% Excludes: Model retrieval (previous), final selection (next)

\begin{figure}[H]
    \centering
    \adjustbox{width=0.7\textwidth,height=0.45\textheight,center,keepaspectratio}{\includegraphics{figures/prediction-option-selection.pdf}}
    \caption{Adaptation option evaluation and selection using learned Q-values (Stage 7). 
    This focused view shows decision-making only. Complete decision workflow in Appendix A.2.}
    \label{fig:option-selection}
\end{figure}

\subsection{Pure Inference at Runtime}

\textbf{Process boundary}: This planning process is \textbf{pure inference}---it uses a pre-trained model to evaluate options but performs no learning, no replay memory sampling, and no policy updates.

% ============================================================================
% SECTION 5.7: Adaptation Option Verification
% ============================================================================

\section{Adaptation Option Verification}
\label{sec:verification}

\subsection{Runtime Model Parameterization}

The predicted adaptation option undergoes verification using runtime models parameterized with current state and candidate option.

% NOTE: Decomposed figure showing only verification trigger and model setup
% Figure to be created: verification-trigger.pdf
% Shows: Predicted option → Runtime model retrieval → Model parameterization
% Excludes: SMC execution, quality evaluation (next figures)

\begin{figure}[H]
    \centering
    \adjustbox{width=0.7\textwidth,height=0.45\textheight,center,keepaspectratio}{\includegraphics{figures/verification-trigger.pdf}}
    \caption{Verification initiation and runtime model parameterization. 
    Complete verification workflow in Appendix B.}
    \label{fig:verification-trigger}
\end{figure}

\subsection{Statistical Model Checking}

Verification employs statistical model checking (UPPAAL-SMC) to estimate goal satisfaction probabilities under uncertainty.

% NOTE: Decomposed figure showing only SMC execution
% Figure to be created: verification-smc.pdf
% Shows: Parameterized model → UPPAAL-SMC → Simulation results
% Excludes: Model setup (previous), quality extraction (next)

\begin{figure}[H]
    \centering
    \adjustbox{width=0.7\textwidth,height=0.45\textheight,center,keepaspectratio}{\includegraphics{figures/verification-smc.pdf}}
    \caption{Statistical model checking execution using UPPAAL-SMC. 
    Complete verification workflow in Appendix B.}
    \label{fig:verification-smc}
\end{figure}

\subsection{Quality Evaluation and Decision Justification}

Verification generates quality estimates that justify adaptation decisions. During runtime, verification provides \textbf{decision justification}. During learning, it generates \textbf{reward signals} as a by-product.

\begin{figure}[H]
    \centering
    \adjustbox{width=0.7\textwidth,height=0.5\textheight,center,keepaspectratio}{\includegraphics{figures/verification-quality-evaluation.pdf}}
    \caption{Quality evaluation and feedback generation from verification results. 
    Produces decision justification (runtime) and learning rewards (asynchronous).}
    \label{fig:quality-evaluation}
\end{figure}

For complete verification workflow including all steps, see Appendix B.

% ============================================================================
% SECTION 5.8: Execution and System Reconfiguration
% ============================================================================

\section{Execution and System Reconfiguration}
\label{sec:execution}

\subsection{Adaptation Option Retrieval}

The Executor retrieves detailed specifications for the selected adaptation option from the Knowledge repository (Stage 8).

\subsection{Adaptation Application}

The Executor translates the adaptation specification into concrete system modifications (Stage 9), applying configuration changes to the managed system.

\begin{figure}[H]
    \centering
    \adjustbox{width=0.6\textwidth,height=0.45\textheight,center,keepaspectratio}{\includegraphics{figures/mape-k-execution-phase.pdf}}
    \caption{Execution and system reconfiguration (Stages 8--9). 
    This focused view shows adaptation enactment only.}
    \label{fig:execution}
\end{figure}

\subsection{Loop Closure}

After adaptation execution, monitoring resumes (Stage 1), closing the adaptation loop and ensuring continuous operation.

% ============================================================================
% SECTION 5.9: RS-DRL Feedback and Learning Overview
% ============================================================================

\section{RS-DRL Feedback and Learning Overview}
\label{sec:feedback}

\subsection{Role of Feedback}

Feedback serves as the bridge between runtime adaptation and learning, enabling the system to improve policies based on adaptation outcomes.

\subsection{Reward Generation from Verification}

Reward signals are generated from verification outcomes, capturing quality-of-service metrics and goal satisfaction indicators. These rewards guide the learning process toward policies that achieve adaptation goals.

\subsection{Separation Between Feedback Generation and Learning}

The architecture maintains clear separation between:
\begin{itemize}
    \item \textbf{Feedback generation}: Transforms outcomes into learning signals
    \item \textbf{Policy learning}: Consumes signals to update models (Section~\ref{sec:learning})
\end{itemize}

% ============================================================================
% SECTION 5.10: Asynchronous RS-DRL Learning and Policy Update
% ============================================================================

\section{Asynchronous RS-DRL Learning and Policy Update}
\label{sec:learning}

This section provides an overview of the RS-DRL learning process using decomposed figures. The complete learning workflow is provided in Appendix C.

\subsection{Learning Lifecycle Overview}

The learning lifecycle operates asynchronously with runtime adaptation:
\begin{enumerate}
    \item Experience collection
    \item Experience storage in replay memory
    \item Model training with sampled experiences
    \item Model storage for runtime use (Stage 3)
\end{enumerate}

\subsection{Experience Collection}

During training episodes, the system collects experiences through environment interaction, generating experience tuples $\langle s_t, a_t, r_t, s_{t+1} \rangle$.

% NOTE: Decomposed figure showing only experience collection
% Figure to be created: training-experience-collection.pdf
% Shows: Episode initialization → Action selection → Environment interaction → Experience tuple
% Excludes: Replay memory, training (next figures)

\begin{figure}[H]
    \centering
    \adjustbox{width=0.75\textwidth,height=0.5\textheight,center,keepaspectratio}{\includegraphics{figures/training-episode-initialization.pdf}}
    \caption{Experience collection during training episodes. 
    This focused view shows episode initialization and interaction only. 
    Complete learning workflow in Appendix C.}
    \label{fig:experience-collection}
\end{figure}

\subsection{Novel Reward-Shaping Mechanism}

The \textbf{reward-shaping mechanism} represents the key innovation of the RS-DRL approach. This mechanism:

\begin{itemize}
    \item Selectively reshapes failed experiences
    \item Treats certain failures as optimistic successes
    \item Accelerates learning in heterogeneous environments
    \item Promotes policy generalization
\end{itemize}

The reward-shaping process operates during model training, transforming the experience distribution to facilitate faster convergence.

% NOTE: Decomposed figure showing only reward reshaping
% Figure to be created: training-reward-reshaping.pdf
% Shows: Minibatch sampling → Failure detection → Reward reshaping → Reshaped batch
% Excludes: Q-network training (next figure)

\begin{figure}[H]
    \centering
    \adjustbox{width=0.75\textwidth,height=0.5\textheight,center,keepaspectratio}{\includegraphics{figures/training-experience-storage.pdf}}
    \caption{Experience replay and novel reward-shaping mechanism. 
    This focused view highlights the key innovation. Complete learning workflow in Appendix C.}
    \label{fig:reward-shaping}
\end{figure}

\subsection{Model Training and Storage}

Model training updates the deep Q-network parameters through gradient descent. The updated model is stored in the Knowledge repository (Stage 3), making it available for future runtime adaptation (Stage 6).

% NOTE: Decomposed figure showing only Q-network training
% Figure to be created: training-q-network.pdf
% Shows: Reshaped batch → TD target computation → Loss computation → Gradient descent → Updated model
% Excludes: Experience collection, reward shaping (previous figures)

\begin{figure}[H]
    \centering
    \adjustbox{width=0.75\textwidth,height=0.5\textheight,center,keepaspectratio}{\includegraphics{figures/training-q-network.pdf}}
    \caption{Q-network training and model update. 
    This focused view shows policy improvement only. Complete learning workflow in Appendix C.}
    \label{fig:q-network-training}
\end{figure}

For the complete end-to-end learning workflow showing all steps integrated, see Appendix C.

% ============================================================================
% SECTION 5.11: Knowledge Management and Process Coordination
% ============================================================================

\section{Knowledge Management and Process Coordination}
\label{sec:knowledge}

\subsection{Knowledge as Shared System Memory}

The Knowledge repository serves as shared system memory, storing trained models, adaptation options, and state history.

\subsection{Read--Write Coordination Contracts}

Knowledge coordination occurs through well-defined contracts:

\begin{itemize}
    \item \textbf{When monitoring completes} (Stage 1): State vectors stored
    \item \textbf{When learning completes} (Stage 3): Trained models written
    \item \textbf{When planning begins} (Stage 6): Trained models read
    \item \textbf{When execution begins} (Stage 8): Adaptation specifications read
\end{itemize}

\subsection{Coordination of Heterogeneous Processes}

The Knowledge repository enables processes with different temporal characteristics to synchronize effectively without tight coupling.

% ============================================================================
% SECTION 5.12: Chapter Summary
% ============================================================================

\section{Chapter Summary}
\label{sec:summary}

\subsection{Hierarchical Process Decomposition}

This chapter presented the architecture and process model through progressive refinement:

\begin{enumerate}
    \item System-level adaptation loop (black-box view)
    \item MAPE-K conceptual mapping (classification lens)
    \item Runtime monitoring and state construction (perception)
    \item Analysis and triggering (control initiation)
    \item Planning and prediction (decision intelligence)
    \item Verification (uncertainty evaluation)
    \item Execution (action)
    \item Learning overview (policy improvement)
    \item Knowledge management (coordination)
\end{enumerate}

\subsection{Architectural Contributions}

Key architectural contributions include:

\begin{itemize}
    \item Temporal separation of runtime decision-making and asynchronous learning
    \item Process-based traceability through nine-stage flow
    \item Verification as justification for learning-based decisions
    \item Knowledge as coordination mechanism for heterogeneous processes
\end{itemize}

\subsection{Figure Organization Principle}

This chapter employed \textbf{decomposed figures} aligned with section structure, where each figure focuses on one conceptual responsibility. This approach:

\begin{itemize}
    \item Reduces cognitive load for examiners
    \item Aligns visual artifacts with textual modularity
    \item Preserves complete workflows in appendices
    \item Signals PhD-level presentation judgment
\end{itemize}

\subsection{Transition to Evaluation}

Having established the architectural foundation, the thesis proceeds to experimental evaluation (Chapter 6), demonstrating how the architectural design decisions enable effective adaptation under uncertainty.

\end{document}
